\section*{Procedimentos}

Nesta secção devem apresentar-se os métodos utilizados para o estudo aerodinâmico do perfil. Neste trabalho, os métodos que os autores devem empregar são: método dos perfis delgados, método dos painéis e simulação CFD.

% Quebra de página explícita presente no ficheiro DOCX original.

\newpage
\enlargethispage{-35pt} % Ajuste da paginação da página 2 da template original

\subsection*{Método dos perfis delgados}

\begin{itemize}[itemsep=3pt]

	\item Deve incluir-se um diagrama de blocos que permita identificar claramente as etapas a seguir para a aplicação do método dos perfis delgados. As etapas apresentadas neste diagrama devem servir de estrutura para a explicação do método nesta secção. Pode adicionar-se formulação matemática caso se considere necessário para explicar o procedimento seguido.

	\item Deve incluir-se uma figura onde se mostre a geometria analisada e se identifique claramente a linha média ou de curvatura do perfil.

\end{itemize}

\begin{figure}[ht]
	\centering

	\def\blockwidth{45mm}
	\begin{tikzpicture}[
			>=Latex, font=\small,
			block/.style={draw=black!75, rounded corners=1.5pt, fill=black!3,
					align=center, inner sep=3mm, minimum height=23mm, text width=\blockwidth},
			flow/.style={->, semithick, draw=black!80},
			phase/.style={font=\small\bfseries, align=center}
		]
		\node[block] (b1) at (0,0) {\textbf{1. Geometria}\\ Coordenadas do extradorso e intradorso\\ $\longrightarrow (x_u,z_u),(x_l,z_l)$};
		\node[block] (b2) at (5.9,0) {\textbf{2. Interpolação e linha média}\\ Entrada: superfícies, $N$\\ Saída: $z(x)=(z_u+z_l)/2$};
		\node[block] (b3) at (11.8,0) {\textbf{3. Diferenciação numérica}\\ Entrada: $z(x)$\\ Saída: $dz/dx$};
		\node[block] (b4) at (11.8,-3.6) {\textbf{4. Mudança de variável}\\ Entrada: $dz/dx$, $N$\\ Saída: $\left.\frac{dz}{dx}\right|_{x(\theta)}$, $\frac{x}{c}=\frac{1-\cos\theta}{2}$};
		\node[block] (b5) at (5.9,-3.6) {\textbf{5. Coeficientes de Fourier}\\ Entrada: declive em $\theta$, $\alpha$\\ Saída: $A_0(\alpha), A_1$};
		\node[block] (b6) at (0,-3.6) {\textbf{6. Resultados aerodinâmicos}\\ Entrada: $A_0,A_1$\\ Saída: $c_l(\alpha),\alpha_{L=0}, dc_l/d\alpha$};
		\draw[flow] (b1) -- (b2);
		\draw[flow] (b2) -- (b3);
		\draw[flow] (b3) -- (b4);
		\draw[flow] (b4) -- (b5);
		\draw[flow] (b5) -- (b6);
	\end{tikzpicture}

	\caption{Diagrama de Blocos - Método dos perfís delgados}
	\label{fig:block-diag-thin-foils}
\end{figure}

\subsection*{Método dos painéis}

\begin{itemize}[itemsep=3pt]

	\item Deve incluir-se um diagrama de blocos que explique o procedimento seguido para a aplicação do método dos painéis. À semelhança do caso anterior, este diagrama deve servir de base para a explicação do método e, caso seja necessário, podem incluir-se no texto as equações consideradas relevantes para fundamentar a explicação.

\end{itemize}

\subsection*{Simulação CFD}

\begin{itemize}[itemsep=3pt]

	\item Deve incluir-se, tal como nos métodos anteriores, a descrição do \textit{setup} utilizado, que explique o procedimento seguido para a realização da simulação CFD.

\end{itemize}
